\begin{tuluyen}
Một học sinh thực hiện phản ứng giữa zinc (Zn) với $V$ lít dung dịch hydrochloric acid (HCl) nồng độ $a$ M (lấy dư) và thu thể tích khí thoát ra theo thời gian. Học sinh đó lặp lại thí nghiệm ba lần và thu được kết quả ở bảng sau:
\begin{center}
\begin{tabular}{c|ccc}
\toprule
Thời gian (s) & Lần 1 & Lần 2 & Lần 3\\
\midrule
0  & 0   & 0   & 0\\
30 & 60  & 59  & 61\\
60 & 90  & 89  & 91\\
90 & 102 & 101 & 103\\
\bottomrule
\end{tabular}
\end{center}
\begin{enumerate}
\item Viết phương trình hóa học xảy ra. Nếu tiến hành thí nghiệm với lượng zinc (Zn) như trên với $V$ lít dung dịch HCl có nồng độ 2$a$ M thì khi kết thúc phản ứng thể tích \ce{H2} có thay đổi không? Giải thích.
\item Tính tốc độ trung bình (mL/s) của phản ứng với các mốc thời gian đã cho ở bảng (từ 0--30s; 30s--60s; 60s--90s). Nhận xét và giải thích sự biến đổi tốc độ phản ứng theo thời gian.
\end{enumerate}
\source{Trích nguyên văn Câu 2.2, đề tuyển sinh chuyên Hóa Điện Biên 2026.}
\dapso{(a) \ce{Zn + 2HCl -> ZnCl2 + H2}; thể tích \ce{H2} cuối cùng không đổi; (b) lần lượt \SI{2.0}{\milli\liter\per\second}, \SI{1.0}{\milli\liter\per\second}, \SI{0.4}{\milli\liter\per\second}; tốc độ giảm dần.}
\begin{loigiaiBT}
Phương trình:
\[
\ce{Zn + 2HCl -> ZnCl2 + H2}.
\]
Do HCl được dùng dư nên lượng \ce{H2} cuối cùng do lượng Zn quyết định. Tăng nồng độ HCl làm phản ứng diễn ra nhanh hơn nhưng không làm tăng lượng \ce{H2} khi Zn ban đầu không đổi.

Tại mỗi thời điểm, lấy trung bình ba lần đo:
\[
\overline V_{30}=\frac{60+59+61}{3}=\SI{60}{\milli\liter},
\]
\[
\overline V_{60}=\frac{90+89+91}{3}=\SI{90}{\milli\liter},
\qquad
\overline V_{90}=\frac{102+101+103}{3}=\SI{102}{\milli\liter}.
\]
Do đó:
\[
\overline v_{0-30}=\frac{60-0}{30}=\SI{2.0}{\milli\liter\per\second},
\]
\[
\overline v_{30-60}=\frac{90-60}{30}=\SI{1.0}{\milli\liter\per\second},
\]
\[
\overline v_{60-90}=\frac{102-90}{30}=\SI{0.4}{\milli\liter\per\second}.
\]
Khi phản ứng diễn ra, nồng độ HCl giảm dần nên tốc độ phản ứng giảm.

\textbf{Vậy,} tốc độ trung bình giảm từ \SI{2.0}{\milli\liter\per\second} xuống \SI{1.0}{\milli\liter\per\second} rồi \SI{0.4}{\milli\liter\per\second}; thể tích \ce{H2} cuối cùng không đổi khi tăng nồng độ HCl nhưng lượng Zn ban đầu giữ nguyên và acid vẫn dư.
\end{loigiaiBT}

\end{tuluyen}
